\section{RAG}

\subsection{1. Fachsicht}

\textbf{RAG} (Retrieval-Augmented Generation) liefert \textbf{Chunks aus indexierten Dokumenten} als zus\textbackslash\{\}"atzlichen Kontext f\textbackslash\{\}"ur Modellanfragen. Implementierung im Paket \textbf{`app/rag/`} mit ChromaDB-Anbindung (\texttt{vector\_store.py}, \texttt{retriever.py}, \texttt{context\_builder.py}, \texttt{rag\_pipeline.py}, \textbackslash\{\}ldots\{\}). Die GUI pflegt Quellen und Indexierung unter \textbf{Operations \$\textbackslash\{\}rightarrow\$ Knowledge}; Schalter wie \texttt{rag\_enabled} und \texttt{rag\_space} liegen in \textbf{`AppSettings`}.

\subsection{2. Rollenmatrix}

\begin{center}
\begin{tabular}{ll}
\hline
Rolle & Nutzung \\
\hline
\textbf{Fachanwender} & Knowledge-Workspace; RAG in Chat aktivieren (wenn UI/Settings es erlauben). \\
\textbf{Admin} & \texttt{chromadb}-Installation, Verzeichnis \texttt{./chroma\_db} (typisch), Embedding-Modell in Ollama. \\
\textbf{Entwickler} & \texttt{app/rag/*.py}, \texttt{app/services/knowledge\_service.py}, \texttt{app/gui/knowledge\_backend.py}. \\
\textbf{Business} & ?Wissensbasierte Antworten\texttt{}. \\
\hline
\end{tabular}
\end{center}

\subsection{3. Prozesssicht}

\begin{verbatim}
User-Query (Chat)
        ?
        v
KnowledgeService / RAG-Pfad (wenn rag_enabled)
        ?
        v
Retriever $\rightarrow$ Vector Store (ChromaDB)
        ?
        v
context_builder / rag_pipeline $\rightarrow$ angereicherter Prompt
        ?
        v
LLM (gleicher Pfad wie Chat)
\end{verbatim}

\textbf{Wichtige Module:} \texttt{service.py}, \texttt{retriever.py}, \texttt{vector\_store.py}, \texttt{context\_builder.py}, \texttt{embedding\_service.py}, \texttt{document\_loader.py}, \texttt{chunker.py}, \texttt{knowledge\_space.py}.

\subsection{4. Interaktionssicht}

\textbf{UI}

\begin{itemize}
\item \texttt{operations\_knowledge} \$\textbackslash\{\}rightarrow\$ \texttt{KnowledgeWorkspace} (\texttt{operations\_screen.py})
\item Backend-Br\textbackslash\{\}"ucke: \texttt{app/gui/knowledge\_backend.py}
\end{itemize}

\textbf{Services}

\begin{itemize}
\item \texttt{app/services/knowledge\_service.py}
\item Registrierter Service-Name \texttt{rag} / \texttt{knowledge\_service} (siehe \texttt{docs/SYSTEM\_MAP.md})
\end{itemize}

\textbf{Settings (`AppSettings`)}

\begin{itemize}
\item \texttt{rag\_enabled}, \texttt{rag\_space}, \texttt{rag\_top\_k}, \texttt{self\_improving\_enabled}, \texttt{chat\_mode} (RAG-bezogen je nach Verwendung im Code)
\end{itemize}

\textbf{Externe Abh\textbackslash\{\}"angigkeit}

\begin{itemize}
\item Paket \texttt{chromadb} in \texttt{requirements.txt}
\item Optional: \texttt{scripts/index\_rag.py} f\textbackslash\{\}"ur Indexierung (referenziert in \texttt{help/troubleshooting/troubleshooting.md})
\end{itemize}

\subsection{5. Fehler- / Eskalationssicht}

\begin{center}
\begin{tabular}{ll}
\hline
Problem & Ursache \\
\hline
Keine Treffer & Leerer Index, falscher \texttt{rag\_space}, Embedding-Modell fehlt. \\
ChromaDB-Fehler & Paket nicht installiert, korruptes Datenverzeichnis. \\
Hohe Latenz & Gro\textbackslash\{\}ss\{\}es \texttt{rag\_top\_k} oder gro\textbackslash\{\}ss\{\}e Sammlung. \\
\hline
\end{tabular}
\end{center}

\subsection{6. Wissenssicht}

\begin{center}
\begin{tabular}{ll}
\hline
Begriff & Ort \\
\hline
\texttt{KnowledgeWorkspace} & \texttt{app/gui/domains/operations/knowledge/knowledge\_workspace.py} \\
\texttt{operations\_knowledge} & Workspace-ID \\
\texttt{vector\_store}, \texttt{retriever} & \texttt{app/rag/} \\
\texttt{rag\_enabled}, \texttt{rag\_space}, \texttt{rag\_top\_k} & \texttt{app/core/config/settings.py} \\
\hline
\end{tabular}
\end{center}

\subsection{7. Perspektivwechsel}

\begin{center}
\begin{tabular}{ll}
\hline
Perspektive & Fokus \\
\hline
\textbf{User} & Knowledge-Workspace, RAG-Toggle in Chat/Settings. \\
\textbf{Admin} & Index-Pipeline, Speicherplatz, Ollama-Embeddings. \\
\textbf{Dev} & Kontrakttests RAG, Failure-Mode-Tests unter \texttt{tests/failure\_modes/}. \\
\hline
\end{tabular}
\end{center}
